\begin{corollary}[Application to Fold: Square-Root Law for Horizon Simulation and \texorpdfstring{$\sqrt\varphi$}{sqrt(phi)} Scaling]\label{cor:dephys-fold-stinespring-sqrt-law}
For fixed resolution \(m\), let \(E_m\) be the fiber-uniform conditional expectation induced by folding, and denote the maximal fiber multiplicity \(D_m=\max_{x\in X_m}d_m(x)\) (Proposition \ref{prop:fold-condexp-index-maxfiber}).
Then any scheme realizing \(E_m\) via finite-dimensional Stinespring dilation must have minimal ancilla dimension \(d_{\min}(m)\) satisfying
\[
d_{\min}(m)\ \ge\ \sqrt{\mathrm{Ind}(E_m)}=\sqrt{D_m}.
\]
Since \(D_m\) satisfies the two-step recurrence \(D_m=D_{m-2}+D_{m-4}\) (Corollary \ref{cor:pom-D-rec}), its exponential growth rate is locked to the Perron root \(r=\sqrt{\varphi}\) of the characteristic equation \(r^4-r^2-1=0\), yielding the order-of-magnitude lower bound
\[
d_{\min}(m)\ \gtrsim\ \varphi^{m/4}.
\]
This inequality binds "the logarithmic scale of maximal information loss \(\log\mathrm{Ind}(E_m)=\log D_m\)" and "the minimal external simulation resource \(\log d_{\min}(m)\)" under a strict square-root relation.
\end{corollary}
